\chapter{Down Syndrome (Trisomy 21)}
\index{Down Syndrome (Trisomy 21)}
\label{sec:Down Syndrome}

\section{Overview}


\begin{itemize}[noitemsep]
  \item Down syndrome is the most common chromosomal disorder, with an incidence of 1 in 700 live births. It is caused by trisomy 21 (94\%), Robertsonian translocation (4\%), or mosaic trisomy 21 (2\%)
  \item Risk increases with maternal age.
\end{itemize}
\section{Causes}
\section{Risk Factors}

\begin{itemize}[noitemsep]
  \item Genetic predisposition and family history
  \item Advancing age
  \item Lifestyle factors (diet, exercise, smoking, alcohol)
  \item Environmental exposures
  \item Underlying medical conditions
  \item Medication use
\end{itemize}
\section{Signs and Symptoms}

\subsection{Common symptoms}
\begin{itemize}[noitemsep]
  \item You may experience flat facial profile, upslanting palpebral fissures, epicanthal folds, Brushfield spots, single transverse palmar crease, hypotonia, joint hyperlaxity, short stature, intellectual disability (mild to moderate), and congenital heart disease (50\%). Associated conditions include hearing loss, hypothyroidism, celiac disease, duodenal atresia, and increased risk of leukemia. Prenatal screening includes first-trimester nuchal translucency and maternal serum markers.

  \item Pain or discomfort in the affected area
  \end{itemize}

\subsection{Emergency warning signs}
\begin{itemize}[noitemsep]
  \item Severe or worsening symptoms of down syndrome (trisomy 21) require immediate medical evaluation
\end{itemize}
\section{Progression and Stages}
The condition may progress through different stages of severity over time. Early stages often involve mild or subtle symptoms that may be overlooked. Moderate stages involve more noticeable symptoms affecting daily function. Advanced stages may involve significant impairment and complications.
\section{Complications}
\begin{itemize}[noitemsep]
  \item Disease progression leading to worsening symptoms
  \item Secondary infections or injuries
  \item Side effects of treatment
  \item Reduced quality of life
  \item Emotional and psychological impact
\end{itemize}
\section{Diagnosis}
\begin{itemize}[noitemsep]
  \item Your doctor confirms the diagnosis with karyotype or FISH. Management requires lifelong multidisciplinary care
  \item Life expectancy has increased to approximately 60 years.
  \item Comprehensive medical history and physical examination
  \item Laboratory tests to assess organ function
  \item Imaging studies to visualize affected structures
  \item Specialized testing depending on the affected system
  \item Consultation with relevant specialists
\end{itemize}
\section{Prevention}
\begin{itemize}[noitemsep]
  \item Maintain a healthy lifestyle with balanced diet and exercise
  \item Avoid known risk factors
  \item Attend regular medical check-ups for early detection
  \item Follow recommended screening guidelines
\end{itemize}
\section{Treatment Options}
Symptomatic management and supportive care Enzyme replacement therapy for certain conditions Gene therapy (emerging for select conditions) Dietary modifications (e.g., phenylketonuria) Regular monitoring for associated complications Physical and occupational therapy Genetic counseling for family planning Participation in clinical trials for new therapies

\begin{itemize}[noitemsep]
  \item Medications to manage symptoms and address underlying mechanisms
  \item Lifestyle modifications and self-care measures
  \item Physical therapy and rehabilitation
  \item Surgical intervention when indicated
  \item Regular monitoring and follow-up care
\end{itemize}
\section{Living with It}
\begin{itemize}[noitemsep]
  \item Follow your treatment plan as prescribed
  \item Maintain regular follow-up with your healthcare provider
  \item Make necessary lifestyle adjustments
  \item Seek support from family, friends, or support groups
  \item Stay informed about your condition
\end{itemize}
